(a) Since $2$ is invertible, symmetric Gaussian elimination diagonalizes the quadratic form. Each nonzero diagonal coefficient has a square root in $k$, so rescaling the corresponding coordinates gives $f=x_0^2+\cdots+x_r^2$.

(b) If $r=0$, then $f=x_0^2$; if $r=1$, choose $a\in k$ with $a^2=-1$ and factor $f=(x_0+ax_1)(x_0-ax_1)$. Conversely, a reducible quadratic is a product of two linear forms, whose associated symmetric matrix has rank at most $2$. Thus it cannot equal a form of rank $r+1\geq3$.

(c) The gradient is $(2x_0,\ldots,2x_r,0,\ldots,0)$. Hence $\operatorname{Sing}Q=Z(x_0,\ldots,x_r)$, a linear subspace of dimension $n-r-1$, empty exactly when $r=n$.

(d) In the complementary $\mathbb{P}^r=Z(x_{r+1},\ldots,x_n)$, the same equation defines a nonsingular quadric $Q'$. Every point of $Q$ outside the axis is a linear combination of a point of $Q'$ and a point of the axis; conversely every such line satisfies $f=0$. Thus $Q$ is the cone over $Q'$ with that axis.
